\section{Introduction}

% Setiap sitasi diambil dari OpenAlex dan kalimatnya dicocokkan dengan abstrak sumbernya (analysis/Skrip/ambil_rujukan.py,
% analysis/Hasil/abstrak_rujukan.json). Klaim tanpa sumber ditandai \todo.

Large language models (LLMs) build on the Transformer architecture \parencite{vaswani2017}. \textcite{brown2020} showed that a single model can handle many language tasks from only a few examples, and the release of open models \parencite{touvron2023} made it practical to run them outside large data centres. Deployment is therefore no longer tied to the cloud. Models now run at the edge of the network, on phones, and inside web browsers; WebLLM, for instance, performs inference entirely in the browser through WebGPU \parencite{ruan2024webllm}.

Research on these models appears to have divided into streams that seldom cite one another. One stream studies LLM-based autonomous agents \parencite{wang2024agents}. A second asks how language models can run on resource-constrained devices and has produced surveys of edge and mobile-edge LLMs \parencite{friha2024,qu2025,zheng2025} and of small language models \parencite{wang2025slm}. A third applies generative AI and LLMs to operations and supply chains, including inventory management and demand forecasting \parencite{jackson2024,nakhost2025,wamba2023}. Small businesses are a setting in which all three streams could meet, because they often have limited cloud budgets and sensitive data \todo{cite evidence on cloud costs and data sensitivity in small businesses}. Whether the literature connects the streams is not yet clear.

Existing reviews cover one stream at a time. The edge-LLM surveys deal with architectures, optimisation, and deployment techniques \parencite{friha2024,qu2025,zheng2025}. The agent survey explains how agents are built and where they are applied \parencite{wang2024agents}. The supply chain studies examine the capabilities and perceived benefits of generative AI in operations \parencite{jackson2024,wamba2023}, and a related review addresses explainable AI \parencite{kosasih2023}. Bibliometric work on LLMs does exist, but it is either broad, with over 5,000 publications up to early 2023 \parencite{fan2024}, or tied to a single domain such as transportation \parencite{hassan2025}. A bibliometric review of agentic AI in small, medium, and micro enterprises reports a sharp rise since 2022 driven by LLMs and multi-agent systems \parencite{olujimi2025}, yet it leaves out the on-device stream. \todo{recheck in Scopus and Web of Science for earlier bibliometric studies that join the three streams; this search was done only in OpenAlex}

A bibliometric analysis suits this problem because it describes the structure of a fast-growing field and shows where its streams connect or stay apart \parencite{donthu2021,zupic2015}. This study maps the journal literature of 2023 to 2026 across the three streams and asks three questions. RQ1: who and what shaped research on LLMs across agents, on-device deployment, and inventory management (citation analysis)? RQ2: what are the foundations and the current research streams (co-citation analysis and bibliographic coupling)? RQ3: which concepts and gaps are emerging (co-word analysis)? \todo{state the theoretical, methodological, practical, and policy contributions once the discussion is final}
